% ============================================================
%  RESEARCH REPORT — standalone, pdflatex-compilable
%  "Forecasting Time-Varying Intermarket Dependencies
%   Between Cryptocurrencies and Conventional Assets
%   Using Machine Learning"
%  Author: Bogdan Babaev  |  2026
% ============================================================
\documentclass[12pt, a4paper]{article}

% --- packages ---
\usepackage[T1]{fontenc}
\usepackage[utf8]{inputenc}
\usepackage[english]{babel}
\usepackage{lmodern}
\usepackage{geometry}
\geometry{margin=2.5cm}
\usepackage{setspace}
\onehalfspacing

\usepackage{amsmath, amssymb}
\usepackage{booktabs}
\usepackage{graphicx}
\usepackage{float}
\usepackage{caption}
\usepackage{subcaption}
\usepackage{hyperref}
\usepackage{xcolor}
\usepackage{enumitem}
\usepackage{microtype}
\usepackage{parskip}
\setlength{\parskip}{6pt}

\hypersetup{
  colorlinks=true,
  linkcolor=blue!60!black,
  urlcolor=blue!60!black,
  citecolor=blue!60!black
}

% --- title ---
\title{%
  \large\textbf{Research Report}\\[6pt]
  \Large Forecasting Time-Varying Intermarket Dependencies\\
  Between Cryptocurrencies and Conventional Assets\\
  Using Machine Learning\\[10pt]
  \normalsize Master's Thesis — Summary Report
}
\author{Bogdan Babaev}
\date{March 2026}

% ============================================================
\begin{document}

\maketitle
\tableofcontents
\newpage

% ============================================================
\section{Introduction and Motivation}
% ============================================================

Financial markets are increasingly interconnected, and the relationship
between cryptocurrency and conventional assets has become one of the most
actively discussed topics in modern quantitative finance.
Unlike traditional assets, Bitcoin and Ethereum trade continuously,
respond rapidly to sentiment shifts, and can simultaneously behave as
a speculative instrument, an inflation hedge, or a macro-risk proxy.
This hybrid nature makes the \emph{dependency} between crypto and
conventional markets time-varying and difficult to model with standard
static tools.

The central question of this thesis is:
\begin{quote}
\textit{Can information derived from cryptocurrency dynamics be used to
forecast time-varying dependency structures with conventional assets,
and can such forecasts be translated into a practically useful
early-warning signal for investors?}
\end{quote}

Answering this question requires a disciplined empirical framework that
avoids the most common pitfalls in financial ML research: in-sample
overfitting, look-ahead bias, and cherry-picked metrics.
The thesis therefore adopts a strict walk-forward out-of-sample
evaluation throughout and compares machine-learning models against
both a classical econometric benchmark (DCC-GARCH) and simple
persistence baselines.

The contribution of the thesis is twofold. First, it provides a
\textbf{reproducible forecasting pipeline} for rolling intermarket
dependencies. Second, it introduces an \textbf{investor signal layer}
that converts dependency forecasts into a stress-warning signal for
conventional markets.

% ============================================================
\section{Data}
% ============================================================

\subsection{Asset universe}

Daily market data were obtained via the \texttt{yfinance} interface.
Bitcoin (\texttt{BTC-USD}) is the base cryptocurrency throughout.
The conventional asset universe consists of six instruments selected
to represent broad equity risk, precious metals, and macro-defensive
positioning.

\begin{table}[H]
  \centering
  \caption{Asset universe.}
  \label{tab:assets}
  \begin{tabular}{lll}
    \toprule
    Ticker        & Description                        & Category \\
    \midrule
    \texttt{BTC-USD}  & Bitcoin                        & Cryptocurrency (base) \\
    \texttt{ETH-USD}  & Ethereum                       & Cryptocurrency (reference) \\
    \texttt{\^{}GSPC} & S\&P 500 index                 & Equity \\
    \texttt{\^{}IXIC} & NASDAQ Composite index          & Equity \\
    \texttt{GLD}      & Gold ETF (SPDR)                & Precious metals \\
    \texttt{SLV}      & Silver ETF (iShares)           & Precious metals \\
    \texttt{UUP}      & U.S.\ Dollar Index ETF proxy   & Macro / FX \\
    \bottomrule
  \end{tabular}
\end{table}

\subsection{Sample period and preprocessing}

The sample spans \textbf{9 November 2017 to 27 February 2026},
yielding 2{,}976 daily observations per asset.
The effective start date is determined by the earliest availability
of ETH-USD on Yahoo Finance.
All price series are converted to daily log returns.
Missing observations due to non-overlapping trading calendars
are handled by alignment and forward-filling of up to two
consecutive non-trading days.

The full dataset therefore contains $7 \times 2{,}975 = 20{,}825$
daily log-return observations, with no remaining missing values
after preprocessing.

\subsection{Dependency targets}

For each of the five Bitcoin–other-asset pairs, a rolling
Pearson correlation is computed over four window lengths
$w \in \{14, 30, 60, 90\}$ trading days, producing
$5 \times 4 = 20$ dependency series in total.

To stabilize the distribution and reduce boundary effects,
each correlation series is passed through the Fisher $z$-transform:
\[
  z_t^{(w)} = \frac{1}{2}\log\!\left(\frac{1 + \rho_t^{(w)}}{1 - \rho_t^{(w)}}\right),
\]
where $\rho_t^{(w)}$ is the rolling correlation at time $t$
with window $w$.
All forecasting models use the Fisher-transformed series as the
prediction target.
The one-step-ahead value $z_{t+1}^{(w)}$ is the quantity to be
forecast from information available at time $t$.

\begin{figure}[H]
  \centering
  \includegraphics[width=0.95\textwidth]{../outputs/figures/eda_rolling_corr_all.png}
  \caption{Rolling correlations between Bitcoin and each of the five conventional
           assets over the full sample (2017--2026), shown for all four window lengths
           ($w = 14, 30, 60, 90$ trading days).
           The shortest window ($w=14$, light blue) is more reactive but noisier;
           longer windows produce smoother, more persistent series.
           Substantial time-variation is visible across all pairs and all window lengths.}
  \label{fig:rolling_corr}
\end{figure}

% ============================================================
\section{Methodology}
% ============================================================

\subsection{Feature engineering}

Predictors are constructed exclusively from information available
at time $t$ or earlier, ensuring no look-ahead contamination.
The feature set covers six interpretable groups:

\begin{table}[H]
  \centering
  \caption{Feature groups used in the dependency-forecasting layer.}
  \label{tab:features}
  \begin{tabular}{ll}
    \toprule
    Group                & Description \\
    \midrule
    Dependency lags      & Recent values of the rolling correlation target \\
    Return lags          & Lagged log-returns of BTC and the paired asset \\
    Rolling means        & Local average return over short sub-windows \\
    Rolling volatility   & Local standard deviation of returns \\
    Spread features      & Signed / absolute gap between BTC and asset behaviour \\
    Cross-asset context  & Additional lags capturing joint dynamics \\
    \bottomrule
  \end{tabular}
\end{table}

Dependency lags are by far the most important predictors because
the rolling-correlation target is smooth and strongly autocorrelated
by construction.
All features are computed inside the training window only;
the test-period window never contributes to feature calculations.

\subsection{Forecasting models}

Eight forecasting approaches are evaluated for each of the 20
dependency experiments, covering the full spectrum from simple
baselines to complex ensemble methods and a classical
econometric benchmark.

\begin{table}[H]
  \centering
  \caption{Models evaluated in the dependency-forecasting layer.}
  \label{tab:models}
  \begin{tabular}{lll}
    \toprule
    Model              & Type                     & Notes \\
    \midrule
    \texttt{Naive\_Last}  & Persistence baseline  & $\hat{z}_{t+1} = z_t$ \\
    \texttt{AR1}          & Autoregressive        & OLS, one lag \\
    \texttt{Ridge}        & Regularised linear    & L2 penalty, standardised features \\
    \texttt{ElasticNet}   & Regularised linear    & L1+L2 penalty \\
    \texttt{RF}           & Tree ensemble         & Random Forest \\
    \texttt{GBM}          & Tree ensemble         & Gradient Boosting \\
    \texttt{XGB\_GPU}     & Tree ensemble         & XGBoost (GPU / CPU fallback) \\
    \texttt{DCC\_GARCH}   & Econometric benchmark & Walk-forward DCC-GARCH(1,1) \\
    \bottomrule
  \end{tabular}
\end{table}

\subsection{Walk-forward evaluation}

All experiments use an \textbf{expanding-window walk-forward design}.
The initial training window contains the first $T_0$ observations
(approximately 60\% of the sample).
At each forecast origin $t \geq T_0$:
\begin{enumerate}[noitemsep]
  \item The model is re-fitted on all observations up to $t$.
  \item A forecast is produced for $t+1$.
  \item The window expands by one day.
\end{enumerate}

This design mirrors real-world forecasting conditions and
ensures that no future information is used in model fitting.
Refitting is performed periodically (every 30 steps) to
balance computational cost and temporal accuracy.

The DCC-GARCH benchmark undergoes the same expanding-window
procedure, making the comparison methodologically fair.
This \emph{leakage-safe} DCC implementation is a key
feature of the thesis: many studies compare ML models
against a full-sample DCC fit, which inadvertently gives
the benchmark access to future data.

\subsection{Evaluation metrics}

Forecasting accuracy is measured by:
\begin{itemize}[noitemsep]
  \item \textbf{RMSE} — root mean squared error (primary ranking metric);
  \item \textbf{MAE} — mean absolute error;
  \item $\boldsymbol{R^2}$ — out-of-sample coefficient of determination.
\end{itemize}

Statistical significance of model differences is assessed with
the \textbf{Diebold--Mariano (DM) test} using a
Newey--West long-run variance estimator.
Two comparisons are made for each experiment:
ML vs.\ DCC-GARCH and ML vs.\ Naive persistence.

\subsection{Investor signal layer}

A second modelling layer reformulates the problem as binary
classification. The target is a \emph{stress-day indicator}
for the conventional asset: a day is labelled as a stress day
if the asset's log return falls below a configurable
negative-sigma threshold.
Three classifiers are tested: \textbf{Logistic Regression},
\textbf{Random Forest (Cls)}, and \textbf{GBM (Cls)}.
Signal performance is measured by Balanced Accuracy,
F1 score for the downside class (F1$_{\text{down}}$),
AUC, and exit rate.

% ============================================================
\section{Results}
% ============================================================

\subsection{Average forecasting performance}

Table~\ref{tab:ranking} presents average out-of-sample performance
across all 20 dependency experiments (5 pairs $\times$ 4 windows).
Models are ranked by average RMSE.

\begin{table}[H]
  \centering
  \caption{Average out-of-sample performance across all 20 dependency experiments.}
  \label{tab:ranking}
  \begin{tabular}{lrrr}
    \toprule
    Model            & Avg.\ RMSE & Avg.\ MAE & Avg.\ $R^2$ \\
    \midrule
    \textbf{AR1}     & \textbf{0.0667} & \textbf{0.0405} & \textbf{0.9421} \\
    Naive\_Last      & 0.0674 & 0.0396 & 0.9409 \\
    Ridge            & 0.0916 & 0.0637 & 0.8914 \\
    ElasticNet       & 0.0923 & 0.0644 & 0.8898 \\
    RF               & 0.1013 & 0.0711 & 0.8687 \\
    GBM              & 0.1038 & 0.0734 & 0.8622 \\
    XGB\_GPU         & 0.1057 & 0.0758 & 0.8566 \\
    \midrule
    DCC\_GARCH       & 0.2345 & 0.1848 & 0.2788 \\
    \bottomrule
  \end{tabular}
\end{table}

Three observations stand out immediately.

\textbf{First}, the best model overall is \texttt{AR1}
(RMSE = 0.0667, $R^2$ = 0.9421).
The naive persistence baseline \texttt{Naive\_Last} is nearly
identical, trailing by only 0.0007 RMSE.
This is not a failure of the framework — it is an empirical
finding about the nature of the target.
Rolling correlation is a smooth, highly autocorrelated series;
the primary source of one-step-ahead predictability is therefore
the target's own recent history.

\textbf{Second}, among the ML models, regularised linear methods
(Ridge and ElasticNet) outperform tree ensembles.
Ridge achieves the best ML RMSE (0.0916), while XGBoost
with GPU acceleration (0.1057) ranks last among ML models.
This pattern is consistent with the high persistence of the target:
ridge shrinkage effectively captures the dominant autoregressive
structure, while boosted trees tend to overfit minor noise.

\textbf{Third}, the DCC-GARCH benchmark is substantially weaker
than all other approaches (RMSE = 0.2345, $R^2$ = 0.2788).
This gap is large and consistent across experiments.

\subsection{Performance across window lengths}

Table~\ref{tab:bywindow} shows the best-model RMSE for each
window length, averaged across the five asset pairs.

\begin{table}[H]
  \centering
  \caption{Best-model RMSE by rolling window length (average across 5 pairs).
           All best models are AR1.}
  \label{tab:bywindow}
  \begin{tabular}{cc}
    \toprule
    Window $w$ & Best RMSE (AR1) \\
    \midrule
    14 days & 0.1451 \\
    30 days & 0.0658 \\
    60 days & 0.0332 \\
    90 days & 0.0228 \\
    \bottomrule
  \end{tabular}
\end{table}

Longer windows produce smoother, more predictable targets;
RMSE drops monotonically from 0.1451 at $w=14$ to 0.0228
at $w=90$.
This confirms that the forecasting difficulty is largely
governed by the smoothness of the target rather than by
the choice of model.

\subsection{Representative forecast visualisation}

Figure~\ref{fig:forecast_sp500} illustrates the walk-forward
forecast for the Bitcoin–S\&P 500 pair at $w=30$, the
experiment of highest practical interest.

\begin{figure}[H]
  \centering
  \includegraphics[width=0.95\textwidth]{../outputs/figures/thesis_forecast_corr_BTC-USD_^GSPC_w30_fisher_z.png}
  \caption{Out-of-sample walk-forward forecast of the Bitcoin–S\&P 500
           rolling dependency (Fisher-transformed, $w=30$).
           The best models closely track the target, especially
           in low-volatility regimes.}
  \label{fig:forecast_sp500}
\end{figure}

The figure highlights a characteristic pattern: models
track broad directional trends well (low-frequency evolution),
but sharp regime transitions — for example, the abrupt
dependency spike during the COVID-19 market crash in
March 2020 — remain difficult for all models.

\subsection{Diebold--Mariano tests}

The Diebold--Mariano framework is used to rigorously test
whether differences in forecast accuracy are statistically
significant.
Two sets of comparisons are performed:

\begin{enumerate}[noitemsep]
  \item \textbf{Ridge vs.\ DCC-GARCH}:
    Ridge is the strongest ML model and the natural
    representative for the ML class.
  \item \textbf{Ridge vs.\ Naive\_Last}:
    Tests whether ML adds value beyond simple persistence.
\end{enumerate}

\textbf{Results — Ridge vs.\ DCC-GARCH}:
All 24 DM tests reject the null of equal predictive accuracy
at the $p < 0.001$ level.
DM statistics range from 20.69 to 33.07, confirming that
Ridge \emph{systematically} outperforms DCC-GARCH across
all dependency experiments.

\textbf{Results — Ridge vs.\ Naive\_Last}:
All 24 DM statistics are \emph{negative} (range: $-12.72$
to $-5.78$), meaning the null of equal accuracy is rejected
in the \emph{opposite} direction — Naive persistence
significantly outperforms Ridge in all experiments.

Figure~\ref{fig:dm_heatmap} summarises the DM comparison
against DCC-GARCH.

\begin{figure}[H]
  \centering
  \includegraphics[width=0.82\textwidth]{../outputs/figures/dm_heatmap_Ridge_vs_DCC.png}
  \caption{Diebold--Mariano test results: Ridge vs.\ DCC-GARCH
           across all 20 dependency experiments.
           Positive DM statistics (warm colours) indicate Ridge wins.
           All tests are significant at $p < 0.001$.}
  \label{fig:dm_heatmap}
\end{figure}

These results support a structured three-level interpretation:
\begin{enumerate}[noitemsep]
  \item \textbf{Is dependency forecastable?} Yes — all models
    achieve high out-of-sample $R^2$ (AR1: 0.942).
  \item \textbf{Does ML beat DCC-GARCH?} Yes — Ridge
    significantly outperforms DCC in all 24 tests.
  \item \textbf{Does ML beat persistence?} No — Naive and AR1
    significantly outperform all ML models on average.
\end{enumerate}

\subsection{Temporal error dynamics}

Forecast accuracy is not constant through time.
Figure~\ref{fig:rolling_rmse} shows the 60-day rolling RMSE
for the Bitcoin--S\&P~500 pair, comparing Ridge (the best ML model),
DCC-GARCH, and the naive persistence baseline.

\begin{figure}[H]
  \centering
  \includegraphics[width=0.92\textwidth]{../outputs/figures/rolling_rmse_corr_BTC-USD_^GSPC_w30_fisher_z.png}
  \caption{60-day rolling RMSE for the Bitcoin--S\&P~500 pair ($w=30$):
           Ridge (best ML model, green) vs DCC-GARCH (blue) vs Naive persistence (dashed).
           Ridge consistently outperforms DCC; the Naive baseline remains the tightest
           overall, confirming that persistence drives most predictability.
           Error spikes cluster around known market stress episodes.}
  \label{fig:rolling_rmse}
\end{figure}

Error spikes cluster around known market stress periods
(2020 COVID crash, 2022 crypto bear market, 2022–2023
rate-hike cycle).
During these episodes all models deteriorate, but the
deterioration is proportionally smaller for the simpler models,
which underreact to abrupt changes without accumulating
additional noise from more complex fitting procedures.

\subsection{Investor signal layer}

Table~\ref{tab:signal} reports average classification results
across all dependency experiments and window lengths.

\begin{table}[H]
  \centering
  \caption{Average investor-signal performance by classifier.}
  \label{tab:signal}
  \begin{tabular}{lrrrr}
    \toprule
    Signal model & Bal.\ Accuracy & F1$_{\text{down}}$ & AUC & Exit Rate \\
    \midrule
    \textbf{Logit}  & \textbf{0.5058} & \textbf{0.2087} & 0.5080 & 0.292 \\
    GBM\_Cls        & 0.5008          & 0.0334          & \textbf{0.5201} & 0.018 \\
    RF\_Cls         & 0.5003          & 0.0297          & 0.5195 & 0.017 \\
    \bottomrule
  \end{tabular}
\end{table}

Logistic regression emerges as the most balanced classifier.
It achieves a Balanced Accuracy of 0.506, meaningfully above
the 0.50 random baseline, with an F1$_{\text{down}}$ of 0.209.
Tree-based classifiers achieve higher AUC but collapse into
near-zero exit rates, effectively failing to flag any
stress days in practice.

The best individual configuration is BTC–S\&P 500, $w=30$,
Logistic Regression:
\[
  \text{Balanced Accuracy} = 0.541,\quad
  \text{F1}_{\text{down}} = 0.252,\quad
  \text{AUC} = 0.538,\quad
  \text{Exit Rate} = 33\%.
\]

Figure~\ref{fig:signal_sp500} shows the generated signal
alongside the S\&P 500 return series.

\begin{figure}[H]
  \centering
  \includegraphics[width=0.95\textwidth]{../outputs/figures/signal_corr_BTC-USD_^GSPC_w30.png}
  \caption{Investor stress-warning signal for the S\&P 500
           ($w=30$, Logit classifier). Red markers indicate
           flagged elevated-risk days.}
  \label{fig:signal_sp500}
\end{figure}

The signal is not strong enough to serve as a standalone
trading rule.
It is best interpreted as a \textbf{risk-management overlay}:
a tool that raises caution flags on roughly 30\% of days,
with moderate success at capturing genuine downside episodes.

% ============================================================
\section{Key Findings and Discussion}
% ============================================================

\subsection{The persistence result is the main finding}

The dominant empirical result is that the rolling dependency
between Bitcoin and conventional assets is forecastable,
but that the forecast is driven almost entirely by
the \textbf{target's own inertia}.
This finding is not a limitation — it is substantively meaningful.

It implies that intermarket dependencies in crypto-conventional
pairs move gradually and persistently.
A portfolio manager who tracks the 30-day rolling BTC–S\&P 500
correlation obtains approximately 94\% of the next day's
variance explained simply by observing today's value.
Complex models cannot easily improve on this because the
signal-to-noise ratio in the \emph{incremental} component is low.

\subsection{Where machine learning adds genuine value}

Machine learning does not beat simple baselines, but it
beats DCC-GARCH \emph{consistently and significantly}
(DM statistics 20–33, all $p < 0.001$).
This has a clear explanation: DCC-GARCH is designed to
model the full conditional covariance structure of returns.
When evaluated purely as a point forecaster of a smoothed
rolling-correlation target, its parametric structure
is less adaptive than even a regularised linear model
that has direct access to recent dependency lags.

This finding validates the design choice of treating rolling
correlation as a direct ML forecasting target, rather than
deriving it from a parametric GARCH system.

\subsection{Practical interpretation of the signal layer}

The signal layer deliberately avoids overstatement.
Results show modest but statistically non-trivial performance
(Balanced Accuracy 0.506–0.541 for Logit).
From a practical standpoint, the signal's value lies in
\emph{risk reduction} rather than \emph{alpha generation}.
A cautious investor who reduces equity exposure on the
$\approx 33\%$ of days flagged by the logistic signal
for the BTC–S\&P 500 pair would skip a disproportionate
share of severe downside days, at the cost of also missing
some positive days.

% ============================================================
\section{Limitations and Future Work}
% ============================================================

\subsection{Limitations}

\begin{itemize}
  \item \textbf{Daily frequency only}: intraday spillovers
    and very short-lived contagion are not captured.
  \item \textbf{One-step-ahead horizon}: multi-day and weekly
    forecasting may reveal different model rankings.
  \item \textbf{Rolling correlation as target}: alternative
    dependence measures (copulas, tail dependence, realised
    covariance) could uncover additional structure.
  \item \textbf{No portfolio backtest}: the signal layer is
    evaluated statistically but not translated into full
    portfolio-level profit/loss accounting with transaction costs.
  \item \textbf{Live data dependency}: results reflect
    market conditions through February 2026; a later
    replication may capture structurally different
    crypto–equity regimes.
\end{itemize}

\subsection{Directions for future research}

\begin{enumerate}[noitemsep]
  \item Extend to multi-day and weekly forecast horizons.
  \item Enrich features with macro variables (VIX, credit spreads,
    on-chain Bitcoin metrics).
  \item Apply richer dependence targets: copula-based
    or realised covariance measures.
  \item Integrate the signal layer into a portfolio
    optimisation framework with turnover constraints.
  \item Evaluate out-of-sample stability across distinct
    crypto market regimes (bull 2020–21, bear 2022,
    recovery 2023–24, new ATH cycle 2024–26).
\end{enumerate}

% ============================================================
\section{Conclusion}
% ============================================================

This report summarises a master's thesis that develops and
evaluates a reproducible two-layer machine-learning pipeline
for forecasting time-varying dependency between Bitcoin
and five conventional asset classes.

The empirical work covers 2{,}976 daily observations
(November 2017 – February 2026), 20 dependency experiments
(5 pairs $\times$ 4 windows), 8 forecasting models, 48 DM tests,
and an investor signal layer with 3 classifiers.

The principal findings are:

\begin{enumerate}
  \item \textbf{Dependency is forecastable} in an out-of-sample
    walk-forward setting across all pairs and windows
    ($R^2 \approx 0.84$–$0.97$ for the best models).
  \item \textbf{Persistence dominates}: AR1 and Naive baselines
    rank first because the smoothed rolling-correlation target
    is strongly autocorrelated.
  \item \textbf{Machine learning beats DCC-GARCH}:
    Ridge significantly outperforms the leakage-safe DCC benchmark
    in all 24 DM tests ($p < 0.001$, DM statistics 20–33).
  \item \textbf{The investor signal is moderate but meaningful}:
    Logistic regression achieves Balanced Accuracy up to 0.54
    and F1$_{\text{down}}$ up to 0.25 for BTC–S\&P 500,
    supporting its use as a risk-management overlay.
\end{enumerate}

These results are internally consistent, methodologically grounded,
and practically interpretable.
The thesis does not claim that cryptocurrency data can predict
conventional markets with high accuracy.
It does show that cryptocurrency dynamics carry structured
information about changing cross-market dependencies,
that this structure is forecastable, and that such forecasts
can inform an investor-facing risk warning mechanism.

% ============================================================
\end{document}
